\begin{question}

  \begin{questionpart}
    Suppose that $f(A)$ is a function of a Hermitian operator $A$ with
    the property $A\ket{a'} = a'\ket{a'}$.  Evaluate
    $\mel{b''}{f(A)}{b'}$ when the transformation matrix from the $a'$
    basis to the $b'$ basis is known.
  \end{questionpart}

  \begin{questionpart}
    Using the continuum analogue of the result obtained in (a),
    evaluate
    \begin{equation*}
      \mel{p''}{F(r)}{p'}
    \end{equation*}
    Simplify your expression as far as you can.  Note that $r$ is
    $\sqrt{x^2 + y^2 + z^2}$, where $x$, $y$, and $z$ are operators.
  \end{questionpart}
\end{question}

\begin{purpose}
  This one seems to be intended to get you to think about clever ways
  to employ unitary operators.  It's actually fun.  I'm feeling the
  get-there-itis\footnote{Always thought this was a dumb name.} pretty
  strongly right now but this one is actually kinda fun.  It's a good
  exploration of a tool and the nuts-and-bolts mechanics of switching
  to/from discrete/continuous spectra.
\end{purpose}

\begin{answer}

  \begin{answerpart}{Unitarius Containus}

    \begin{equation}
      \begin{split}
        \mel{b_i}{f(A)}{b_j}
        &= \bra{a_i}U^\dagger f(A)U\ket{a_j}\\
      \end{split}
    \end{equation}

  Given what we're given in the problem that's pretty much where it
  ends.  We don't know anything about $f(A)$.  Take a look at problem
  \ref{1.6:theQuestion} equation \ref{1.6:kittens}.  Buuut, just like
  we did back in \ref{1.6:theQuestion}, we're going to assume that it
  can be written as a Taylor series.  If that's the case, then we can
  pull in equation \ref{q:taylor} and show that

  \begin{equation}
    \inlabel{mainPoint}
    \begin{alignedat}{2}
      &&f(A)
      &\equiv \sum_{k=0}^{\infty} \frac{1}{k!} A^k
      \\
      \implies &&
      \mel{b_i}{f(A)}{b_j}
      &= \bra{b_i}f(A)\ket{b_j}
      \\
      &&&=
      \sum_{k=0}^{\infty} \frac{1}{k!} \bra{b_i}A^k\ket{b_j}
      \\
      &&&=
      \sum_{k=0}^{\infty} \frac{1}{k!} \bra{a_i} U^\dagger A^k U \ket{a_j}
      \\
    \end{alignedat}
  \end{equation}

  So...what to do what to do what to do about that inner part...

  \begin{equation}
    \begin{alignedat}{2}
      &&A\ket{b_j}
      &= \left(\sum_i a_i\dyad{a_i}{a_i}\right)\ket{b_j}\\
      &&&= \sum_i a_i\braket{a_i}{b_j}\ket{a_i}\\
    \end{alignedat}
  \end{equation}

  So, we know that there exists some unitary operator $U$ such that
  $U\ket{a_i} = \ket{b_i}$.\footnote{Yeah, I know, earlier we said
  that $U$ would turn $\ket{b_i}$ into $\ket{a_i}$, now I'm saying
  it's the other way around, get off my back already.}  Those whacky
  unitarians...always turning one ket into another.\footnote{You
  should see the Unitarian where $U_j\ket{\text{water}} =
  \ket{\text{wine}}$.}  That operator will look like this

  \begin{equation}
    U \dot{=}
    \begin{pmatrix}
      \braket{a_1}{b_1} & \braket{a_2}{b_2} & \ldots\\
      \braket{a_2}{b_1} & \braket{a_2}{b_2} & \ldots\\
      \vdots & \vdots & \ddots\\
    \end{pmatrix}
  \end{equation}

  If we pull in those constants above, let's see what we get
  
  \begin{equation}
    \begin{alignedat}{2}
      &&\mel{b_i}{A}{b_j}
      &= \bra{b_i}\left(\sum_k a_k\dyad{a_k}{a_k}\right)\ket{b_j}\\
      &&&= \bra{b_i}\sum_k a_k\braket{a_k}{b_j}\ket{a_k}\\
      &&&= \sum_k a_kU_{kj}\braket{b_i}{a_k}\\
      &&&= \sum_k a_kU_{kj}U_{ki}^{*}\\
    \end{alignedat}
  \end{equation}

  If we applied A more than once, we'd find that
  
  \begin{equation}
    \begin{alignedat}{2}
      &&\mel{b_i}{A^N}{b_j}
      &= \bra{b_i}A^{N-1}\left(\sum_k a_k\dyad{a_k}{a_k}\right)\ket{b_j}\\
      &&&= \bra{b_i}A^{N-1}\left(\sum_k a_kU_{kj}\ket{a_k}\right)\\
      &&&= \bra{b_i}\left(\sum_k a_k^N U_{kj}\ket{a_k}\right)\\
      &&&= \sum_k a_k^N U_{kj}\braket{b_i}{a_k}\\
      &&&= \sum_k a_k^N U_{kj}U_{ki}^{*}\\
    \end{alignedat}
  \end{equation}

  Armed with this expression, we revisit equation \inref{mainPoint}.
  We're going to make a slight change of counter variable from $k$ to
  $n$

  \begin{equation}
    \begin{alignedat}{2}
      &&\mel{b_i}{f(A)}{b_j}
      &=
      \sum_{n=0}^{\infty} \frac{1}{n!} \bra{b_i} A^n \ket{b_j}
      \\
      &&&=
      \sum_{n=0}^{\infty} \frac{1}{n!} \sum_k a_k^n U_{kj}U_{ki}^{*}
      \\
      &&&=
      \sum_k U_{kj}U_{ki}^{*} \sum_{n=0}^{\infty} \frac{1}{n!}a_k^n
      \\
      &&&=
      \sum_k U_{kj}U_{ki}^{*}f(a_k)\qed
      \\
    \end{alignedat}
  \end{equation}

  That's about as far as I can take it.  Because of the nature of
  having the $\bra{b_i}$ and $\ket{b_j}$ on the outside, it seems like
  you're stuck with a term on the inside for each of those and you'll
  have to sweep the entire basis at least once, so you aren't getting
  rid of the sum.
  \end{answerpart}

  \begin{answerpart}{Unitarius Continuous}

    As we transition from discrete to continuous and from a
    generalized $U$ to a specific $U$ (the momentum operator) we find
    that we can use our above expression as a guide.  We're going to
    start with our original expression and replace one term at a time.
    Because I'm replacing one thing at a time, many of these
    intermediate expressions will be nonsensical when fully evaluated.

    \begin{equation}
      \begin{alignedat}{2}
        &&\underbrace{\mel{p_i}{F(r)}{p_j}}_{\text{New target expression.}}
        &=
        \sum_k U_{kj}U_{ki}^{*}f(a_k)
        \\
      \end{alignedat}
    \end{equation}

    Next up, we'll replace the terms from the unitary transformation
    matrix with their BraKet definitions.

    \begin{equation}
      \begin{alignedat}{2}
        &&\mel{p_i}{F(r)}{p_j}
        &=
        \sum_k \underbrace{\braket{a_k}{b_j}\braket{b_i}{a_k}}_{\text{The terms in the unitary matrix are just BraKets}}f(a_k)
        \\
      \end{alignedat}
    \end{equation}
    
    Time for us to transition from discrete to continuous spectra.
    You have to be a little bit careful here.  Typically definitions
    will employ language like ``This represents the froomba of the
    gworplick when evaluated between a Santa value of $j$ and
    $j+\delta j$ jollies''.  That's typically where you look for
    finding the right scaling factors on your differential (either
    that or you're using the fundamental relation $dX = \sum
    \pdv{}{x_i}X dx_i$).  In this case, I believe it's pretty
    straight-forward.  Don't make the same mistake I did at first, and
    note that the index in the sum is $k$ and that the $k's$ are on
    the $a$'s which is the basis in the function call, not the basis
    of the Bra/Ket giving it a hug.

    \begin{equation}
      \begin{alignedat}{2}
        &&\mel{p_i}{F(r)}{p_j}
        &= \underbrace{\int_{-\infty}^{\infty} \dd{x}}_{\text{Switch to integral}}
        \braket{a_k}{b_j}\braket{b_i}{a_k}f(a_k)
        \\
      \end{alignedat}
    \end{equation}

    Next up, we'll replace our function of $a_k$ with our new function
    of $r$.
    
    \begin{equation}
      \begin{alignedat}{2}
        &&\mel{p_i}{F(r)}{p_j}
        &= \int_{-\infty}^{\infty} dx\ 
        \braket{a_k}{b_j}\braket{b_i}{a_k}\underbrace{F(r)}_{\text{New function}}
        \\
      \end{alignedat}
    \end{equation}

    We're also going to formalize something quickly to avoid any
    confusion.  $F(r)$ is really $F(f(\hat{x}, \hat{y}, \hat{z}))$
    where $f(\hat{x}, \hat{y}, \hat{z}) = \sqrt{\hat{x}^2 + \hat{y}^2
      + \hat{z}^2}$.  We learned in part a that when you have
    functions like that in here, you include your unitary matrix
    elements (the BraKets), loop over all possible eigenstates, and
    you get to replace the operators with eigenvalues.  Next we
    replace the variables in our unitary entries in Dirac form with
    the new variables for this problem.

    \begin{equation}
      \begin{alignedat}{2}
        &&\mel{p_i}{F(r)}{p_j}
        &= \int_{-\infty}^{\infty} dx\ 
        \underbrace{\braket{x}{p_j}\braket{p_i}{x}}_{\text{New Variables}}F(r)
        \\
      \end{alignedat}
    \end{equation}

    Before we move further, it's worth pausing to consider our
    variables.  We have $x$ and $p$ and $r$.  So we have 3 position
    variables in $r$ and only 1 in in our integral.  So we really have
    a generalized 3D position and a generalized 3D momentum.  So let's
    fix our formalism real quick here, just a small change.

    \begin{equation}
      \begin{alignedat}{2}
        &&\mel{p_i}{F(r)}{p_j}
        &= \underbrace{\iiint_{-\infty}^{\infty} \dd[3]{\vb{x}}}_{\text{Triple the pleasure}}
        \braket{\vb{x}}{\vb{p_j}}\braket{\vb{p_i}}{\vb{x}}F(r)
        \\
      \end{alignedat}
    \end{equation}

    So, given that the only unknown is a function of $r$ only, this
    problem is clearly screaming for spherical coordinates.
    
    \begin{equation}
      \begin{alignedat}{2}
        &&\mel{p_i}{F(r)}{p_j}
        &= \int_{0}^{\infty} \int_{0}^{\pi} \int_{0}^{2\pi}
        \braket{\vb{x}}{\vb{p_j}}\braket{\vb{p_i}}{\vb{x}}F(r)
        \underbrace{r^2 \sin(\theta) \dd{r}\dd{\theta}\dd{\varphi}}
        \\
      \end{alignedat}
    \end{equation}

    Now we're down to evaluating those BraKets written using a
    generalized $\vb{x}$ to indicate position.  Before getting into
    those nuts-and-bolts mechanics, let's first pull in a utility
    definition from the book

    \begin{equation}
      \tag{1.282}
      \braket{\vb{x}'}{\vb{p}'} = \frac{1}{(2\pi\hbar)^{3/2}} \exp(\frac{i\vb{p}'\vb{x}'}{\hbar})
    \end{equation}

    Substituting that in, we find that
    
    \begin{equation}
      \begin{alignedat}{2}
        &&\mel{p_i}{F(r)}{p_j}
        &=
        \frac{1}{(2\pi\hbar)^3}
        \int_{0}^{\infty} \int_{0}^{\pi} \int_{0}^{2\pi}
        \underbrace{
          \exp(\frac{i\vb{p_j}\cdot\vb{x}}{\hbar})
          \exp(-\frac{i\vb{p_i}\cdot\vb{x}}{\hbar})}_{
          \text{Explicit expressions for BraKets}}
        F(r)
        r^2 \sin(\theta) \dd{r}\dd{\theta}\dd{\varphi}
        \\
        &&&= \frac{1}{(2\pi\hbar)^3}\int_{0}^{\infty} \int_{0}^{\pi} \int_{0}^{2\pi}
        \exp(\frac{i\vb{x}\cdot(\vb{p_j}-\vb{p_i})}{\hbar})
        F(r)
        r^2 \sin(\theta) \dd{r}\dd{\theta}\dd{\varphi}
        \\
        &&&= \frac{1}{(2\pi\hbar)^3}\int_{0}^{\infty} \int_{0}^{\pi} \int_{0}^{2\pi}
        \exp(\frac{i\vb{x}\cdot\vb{\delta_p}}{\hbar})
        F(r)
        r^2 \sin(\theta) \dd{r}\dd{\theta}\dd{\varphi}
        \\
      \end{alignedat}
    \end{equation}

    The inner product is the defined as $\vb{a}\cdot\vb{b} =
    \abs{a}\abs{b}\cos(\theta)$ where $\theta$ is the angle between
    the two vectors.  We'll rename that angle to be $\gamma$

    \begin{equation}
      \begin{alignedat}{2}
        &&\mel{p_i}{F(r)}{p_j}
        &= \frac{1}{(2\pi\hbar)^3}\int_{0}^{\infty} \int_{0}^{\pi} \int_{0}^{2\pi}
        \exp(\frac{i}{\hbar}r\abs{\delta_p}\cos(\gamma))
        F(r)
        r^2 \sin(\theta) \dd{r}\dd{\theta}\dd{\varphi}
        \\
      \end{alignedat}
    \end{equation}

    Because $F(r)$ is a function \emph{only} of radius, we can choose
    any arbitrary orientation of our $x$, $y$, and $z$ axes that we
    want.  So, if we were to choose our $z$ axis such that it is
    colinear with $+\delta_p$ then the angle between them is exactly
    $\theta$ (irrespective of radius and azimuth).

    \begin{equation}
      \begin{alignedat}{2}
        &&\mel{p_i}{F(r)}{p_j}
        &= \frac{1}{(2\pi\hbar)^3}\int_{0}^{\infty} \int_{0}^{\pi} \int_{0}^{2\pi}
        \exp(\frac{i}{\hbar}r\abs{\delta_p}\cos(\gamma))
        F(r)
        r^2 \sin(\theta) \dd{r}\dd{\theta}\dd{\varphi}
        \\
        &&&= \frac{1}{(2\pi\hbar)^3}\int_{0}^{\infty} \int_{0}^{\pi} \int_{0}^{2\pi}
        \exp(\frac{i}{\hbar}r\abs{\delta_p}\cos(\theta))
        F(r)
        r^2 \sin(\theta) \dd{r}\dd{\theta}\dd{\varphi}
        \\
      \end{alignedat}
    \end{equation}

    We have now removed any azimuthal dependence in the intgral so
    let's pull out our $\varphi$
    
    \begin{equation}
      \begin{alignedat}{2}
        &&\mel{p_i}{F(r)}{p_j}
        &= \frac{1}{(2\pi\hbar)^3}\int_{0}^{\infty} \int_{0}^{\pi} \int_{0}^{2\pi}
        \exp(\frac{i}{\hbar}r\abs{\delta_p}\cos(\theta))
        F(r)
        r^2 \sin(\theta) \dd{r}\dd{\theta}\dd{\varphi}
        \\
        &&&= \frac{1}{(2\pi\hbar)^3}
        \left(\int_{0}^{2\pi} \dd{\varphi}\right)
        \int_{0}^{\infty} \int_{0}^{\pi}
        \exp(\frac{i}{\hbar}r\abs{\delta_p}\cos(\theta))
        F(r)
        r^2 \sin(\theta) \dd{r}\dd{\theta}
        \\
        &&&= 2\pi \frac{1}{(2\pi\hbar)^3}
        \int_{0}^{\infty} \int_{0}^{\pi}
        \exp(\frac{i}{\hbar}r\abs{\delta_p}\cos(\theta))
        F(r)
        r^2 \sin(\theta) \dd{r}\dd{\theta}
        \\
      \end{alignedat}
    \end{equation}

    Shuffling things around a bit and then doing a quick substitution
    of $k \equiv r\abs{\delta_p}/\hbar$ we have
    
    \begin{equation}
      \begin{alignedat}{2}
        &&\mel{p_i}{F(r)}{p_j}
        &= 2\pi \frac{1}{(2\pi\hbar)^3}
        \int_{0}^{\infty} \int_{0}^{\pi}
        \exp(\frac{i}{\hbar}r\abs{\delta_p}\cos(\theta))
        F(r)
        r^2 \sin(\theta) \dd{r}\dd{\theta}
        \\
        &&&= 2\pi \frac{1}{(2\pi\hbar)^3}
        \int_{0}^{\infty}
        r^2 F(r)
        \left(
        \int_{0}^{\pi}
        \exp(ik\cos(\theta))
        \sin(\theta)
        \dd{\theta}
        \right)
        \dd{r}
        \\
      \end{alignedat}
    \end{equation}

    Employing $u$ substitution with $u = \cos(\theta)$ and $du =
    \pdv{}{\theta}\cos(\theta) = -\sin(\theta)$ where the bounds are
    now from $cos(0)=1$ to $cos(pi)=-1$

    \begin{equation}
      \begin{alignedat}{2}
        &&\mel{p_i}{F(r)}{p_j}
        &= 2\pi \frac{1}{(2\pi\hbar)^3}
        \int_{0}^{\infty}
        r^2 F(r)
        \left(
        \int_{0}^{\pi}
        \exp(ik\cos(\theta))
        \sin(\theta)
        \dd{\theta}
        \right)
        \dd{r}
        \\
        &&&= 2\pi \frac{1}{(2\pi\hbar)^3}
        \int_{0}^{\infty}
        r^2 F(r)
        \left(
        \int_{1}^{-1}
        -
        \exp(iku)
        \dd{u}
        \right)
        \dd{r}
        \\
        &&&= 2\pi \frac{1}{(2\pi\hbar)^3}
        \int_{0}^{\infty}
        r^2 F(r)
        \left(
        \int_{-1}^{1}
        \exp(iku)
        \dd{u}
        \right)
        \dd{r}
        \\
        &&&= 2\pi \frac{1}{(2\pi\hbar)^3}
        \int_{0}^{\infty}
        r^2 F(r)
        \eval{[\frac{1}{ik}e^{iku}]}_{-1}^{1}
        \dd{r}
        \\
        &&&= 2\pi \frac{1}{(2\pi\hbar)^3}
        \int_{0}^{\infty}
        r^2 F(r)
        (\frac{1}{ik}(e^{ik} - e^{-ik}))
        \dd{r}
        \\
      \end{alignedat}
    \end{equation}

    Here, we're going to note that $e^{ik} - e^{-ik} = 2isin(k)$ and
    substitute that in.

    \begin{equation}
      \begin{alignedat}{2}
        &&\mel{p_i}{F(r)}{p_j}
        &= 2\pi \frac{1}{(2\pi\hbar)^3}
        \int_{0}^{\infty}
        r^2 F(r)
        (\frac{1}{ik}(e^{ik} - e^{-ik}))
        \dd{r}
        \\
        &&&= 2\pi \frac{1}{(2\pi\hbar)^3}
        \int_{0}^{\infty}
        r^2 F(r)
        \frac{1}{ik}
        2i\sin(k)
        \dd{r}
        \\
        &&&= 2\pi \frac{1}{(2\pi\hbar)^3}
        \int_{0}^{\infty}
        r^2 F(r)
        \frac{\hbar}{ir\abs{\delta_p}}
        2i\sin(k)
        \dd{r}
        \\
        &&&=
        \frac{2\hbar}{\abs{\delta_p}}
        2\pi \frac{1}{(2\pi\hbar)^3}
        \int_{0}^{\infty}
        r F(r)
        \sin(r\frac{\abs{\delta_p}}{\hbar})
        \dd{r}
        \\
        &&&=
        \frac{1}{2\abs{\delta_p}}
        \frac{1}{(\pi\hbar)^2}
        \int_{0}^{\infty}
        r F(r)
        \sin(r\frac{\abs{\delta_p}}{\hbar})
        \dd{r}
        \\
      \end{alignedat}
    \end{equation}

    This is the end of the solution, but not the end of my
    pontifications.  I took issue with my own result, even though it
    corresponds to the answer in the official solutions manual, I
    don't see any mistakes, and my AI sycophant says I'm a brilliant
    ray of sunshine.  I didn't like this notion of a dimensionless
    value in that sin that isn't a trigonometric ratio of something.
    It turns out to be related to the wave number (according to the
    aforementioned AI sycophant).

    It sounds like this is one of those cases where if I push too
    hard, I'm going to end up in the Zee book trying to learn
    \ac{qft}.  With $p = k\hbar \implies k = p/\hbar$ we know that
    we're taking the $\sin(r\abs{\vb{k_2} - \vb{k_1}})$ -- so that's the
    radius times the difference in the wave vectors.  I'd love to
    explore more, but I'm guessing that if I pull too hard on that
    thread, I'm going to end up in chapter 6 (Scattering Theory) which
    (according to a buddy of mine who has already been through two
    semesters of it) is the heart and soul of \ac{qft}.

  \end{answerpart}
\end{answer}
